\chapter{Architecture Decisions}\label{app:adr}

% [#90/AP-B5] Compact ADR collection (Architecture Decision Records), traced from the chapters.
% Documented decisions only; every ADR is backed by a \ref to the supporting thesis section.
This appendix collects the load-bearing design decisions of this thesis as compact
\emph{Architecture Decision Records}\footnote{Architecture Decision Record (ADR): a short, numbered
record of a design decision with context, decision and consequences --- introduced by
Nygard~\cite{nygard2011adr} as a short text file in a format in the manner of an Alexander pattern and
described as a tool format by Kopp, Armbruster and Zimmermann~\cite{kopp2018madr}.} (ADRs): for each
decision the context in which it was made, the decision itself and its most important consequence. The
goal is quick traceability --- \emph{why} the apparatus is built the way the chapters describe it; the
detailed derivation and justification resides in the section of the main text referenced under
\emph{Evidence} and is not repeated here. Where design and implementation diverge, the ADR states the
implementation status at exactly one place as a fourth field. The collection comprises twenty
decisions: the twelve decisions of the conception (ADR-1 to ADR-12) and eight decisions of identity,
interfaces, operation and build (ADR-13 to ADR-20); Table~\ref{tab:adr-overview} gives the overview.

\begin{footnotesize}
\begin{longtable}{@{}>{\raggedright\arraybackslash}p{1.3cm}>{\raggedright\arraybackslash}p{6.0cm}>{\raggedright\arraybackslash}p{3.3cm}>{\raggedright\arraybackslash}p{3.8cm}@{}}
\caption{Overview of the twenty architecture decisions: short title, evidencing section and
implementation status --- \emph{implemented} (implemented and evidenced), \emph{partial}
(implemented, with a declared open item) or \emph{design} (designed only, not implemented); the
justifications are in the records, not in the table}\label{tab:adr-overview}\\
\toprule
ADR & Decision (short title) & Evidence (section) & Implementation status \\
\midrule
\endfirsthead%
\multicolumn{4}{@{}l}{\tablename\ \thetable{} --- continued}\\
\toprule
ADR & Decision (short title) & Evidence (section) & Implementation status \\
\midrule
\endhead%
\midrule
\multicolumn{4}{@{}r}{\footnotesize continued on the next page}\\
\endfoot%
\bottomrule
\endlastfoot%
ADR-1 & C-stable module boundary instead of a C++ module ABI & Sect.~\ref{sec:impl-architecture} & implemented \\
ADR-2 & Compile-time permutation via type lists & Sect.~\ref{sec:axes-impl} & partial (activation of the eighteenth axis) \\
ADR-3 & Three-level model as the one architecture & Sect.~\ref{sec:anatomy-levels} & implemented \\
ADR-4 & Pass-through doctrine and enforced full assignment & Sect.~\ref{sec:axes-impl} & partial (pass-through building block of the node type) \\
ADR-5 & Layer model E4--E0 with frozen contracts & Sect.~\ref{sec:anatomy-layers} & implemented \\
ADR-6 & Explicitly coded non-observation instead of a phantom value & Sect.~\ref{sec:measurement-system} & implemented \\
ADR-7 & Four-repository architecture with orthogonal carrier stages & Sect.~\ref{sec:repos} & implemented \\
ADR-8 & Family-wise adapter fallbacks & Sect.~\ref{sec:axes-impl} & implemented \\
ADR-9 & Measurement gating in three layers and four work modes & Sect.~\ref{sec:impl-pipeline} & partial (replay run compare/release) \\
ADR-10 & \texttt{std::map} contract with conformance gate & Sect.~\ref{sec:measurement-system} & partial (full \texttt{std::map} interface, container hull) \\
ADR-11 & Design-pattern register with evidenced roles & Sect.~\ref{sec:anatomy-patterns} & implemented \\
ADR-12 & Dual measurement regime Talos and root Linux & Sect.~\ref{sec:eval-method} & partial (Talos regime) \\
ADR-13 & Four carrier stages and one stamp contract & Sect.~\ref{sec:stamp-model} & partial (stamp export hybrid, SOTA branch) \\
ADR-14 & Identity cut: only organ axes form the \texttt{binary\_id} & Sect.~\ref{sec:axis-triad} & implemented \\
ADR-15 & Measurement-visitor seam (surface 3) instead of a sidecar & Sect.~\ref{sec:anatomy-levels} & implemented \\
ADR-16 & Store as stock with compile-time keys & Sect.~\ref{sec:lager} & partial (full store build-out) \\
ADR-17 & Optimization level as the user's choice, build world identity-effective & Sect.~\ref{sec:eval-build-scope} & partial (O0/O1 choice, flag pass-through) \\
ADR-18 & One official build route as a GNU front end over CMake & Appendix~\ref{app:code} & partial (framework export) \\
ADR-19 & Four CI cells without skip & Sect.~\ref{sec:impl-contracts} & partial (mini-pipelines per carrier stage) \\
ADR-20 & XML as the only control & Sect.~\ref{sec:impl-pipeline} & partial (creation mode, headless control path, schema drift) \\
\end{longtable}
\end{footnotesize}

\begin{description}

  \item[ADR-1 --- C-stable module boundary instead of a C++ module ABI]\hfill\\
  \textbf{Context:} Each permutation is compiled as a standalone tier binary and loaded at run
  time; C++23 modules structure the translation units but, by themselves, guarantee no portable
  ABI stability~\cite{wg21_modules_abi}.
  \textbf{Decision:} A uniform C module boundary exports exactly one living being per permutation
  --- an \texttt{IAnatomyBase} instance --- via seven mandatory symbols named independently of the
  family: the four primal symbols ABI version, magic, factory (\texttt{comdare\_create\_anatomy()})
  and destroy, the two identity symbols family and genus, whereby the family is not maintained as a
  second datum but derived statically at compile time from the genus, and the stamp-line symbol. The
  identity of names across all families is the stability argument: the loader can determine the
  family via a family-independent symbol without knowing it in advance. Behind it lie two engine
  concepts: the \texttt{CacheEngine} as the library of the axis building blocks and the
  \texttt{ExecutionEngine} as the execution and measurement infrastructure with the root
  \texttt{IExecutionEngine}, under which \texttt{IAnatomyBase} stands; the API facade is
  \texttt{ICacheEngine}; the name \enquote{SearchEngine} denotes the ABI runtime view of the genus
  \texttt{SearchAlgorithm}, not a third link in the chain; the same-named module intended to house it
  separately is declared but not wired in.
  \textbf{Consequence:} ABI stability is established by the C boundary, not by the module system
  (ABI major~9); the loader checks every module on loading, fail-closed, against the mandatory
  symbol set and the identity derivation --- a module with contradictory identity or without a stamp
  symbol is rejected rather than wandering into stamp, store and measurement series as a wrong
  identity.
  \textbf{Evidence:} Sections~\ref{sec:contributions},~\ref{sec:anatomy-patterns}
  and~\ref{sec:impl-architecture}; Figure~\ref{fig:one-architecture}.

  \item[ADR-2 --- Compile-time permutation via type lists instead of runtime adaptation]\hfill\\
  \textbf{Context:} Eighteen organ composition axes already count $2^{17} = 131072$ binary
  identities per system permutation in the canonical measurement catalogue --- one of them, the
  persistence target, pinned to its in-memory representative ---, and the raw variant product space
  over all axis catalogues including their sub-axes lies orders of magnitude above that; the hot
  path must carry neither \texttt{virtual} calls nor runtime switches. The pinning is a budget
  decision (according to the catalogue header about 224~GB and 24~h of complete build instead of
  about 448~GB and 48~h); the variation of all eighteen axes is provided for, and the cardinality of
  a concrete plan is computed by the planner per XML in the director walk, never statically.
  \textbf{Decision:} The permutation runs at compile time via Boost.Mp11 type
  lists\footnote{Boost.Mp11: a C++11 metaprogramming library for the compile-time manipulation of
  data structures that contain types; its list type is \texttt{mp\_list}~\cite{dimov_mp11}.}: per
  axis one full list, filtered via the activation predicate into the enabled list, bound to the
  composition slot, plus concept-checked axis descriptors with the inventory \texttt{variants} and a
  generic registry-to-tree reflection; the golden reference catalogue selects per slot the first
  representatives of the enabled list --- no runtime tag, no \texttt{std::variant}.
  \texttt{std::variant}, \texttt{virtual} and runtime switches are categorically excluded in tier
  and hybrid binaries; a static check enforces this, and the variant-based building-block cluster
  (\texttt{baustein\_variants}, \texttt{resolve\_baustein}) is a test inventory isolated from the
  build, without library consumers.
  \textbf{Consequence:} The axis abstraction remains a zero-cost abstraction; only the eighteen
  organ axes enter the binary identity (ADR-14).
  \textbf{Implementation status:} The inclusion of the eighteenth axis in the reference catalogue is
  not implemented; it requires two steps at once --- catalogue degree two for the persistence target
  and the write-back switch of the persistence axis.
  \textbf{Evidence:} Sections~\ref{sec:contributions},~\ref{sec:anatomy-patterns},~\ref{sec:axes-impl}
  and~\ref{sec:mess-chain}; Figure~\ref{fig:patterns}.

  \item[ADR-3 --- The three-level model as the \emph{one} architecture]\hfill\\
  \textbf{Context:} The anatomy metaphor and the technical engine view threaten to drift apart
  into parallel model trees.
  \textbf{Decision:} There is exactly one hierarchy without parallel trees --- family (level~1,
  outer interface: the probe-dock families Map, Container and Graph as well as the fourth family
  HeuristikAdapter, which has no dock of its own because its binaries report the target genus via
  \texttt{genus()} and dock at the dock of the target) $\to$ genus (level~2, the tier subclass with
  a fixed axis set, alias family subclass; the six genera \texttt{SearchAlgorithm}, \texttt{Set},
  \texttt{Sequence}, \texttt{Adapter}, \texttt{View} and \texttt{FunctionInterfaceReroute}) $\to$
  axis organs (level~3) --- under \texttt{IExecutionEngine} as the common measurable root; a family
  of its own would arise only if the base axis set differed fundamentally.
  \textbf{Consequence:} The genus \texttt{SearchAlgorithm} under the family Map and its ABI runtime
  view (\enquote{SearchEngine}) are two views of the same living being, not a second model.
  \textbf{Evidence:} Sections~\ref{sec:anatomy-model},~\ref{sec:anatomy-levels}
  and~\ref{sec:impl-architecture}; Figures~\ref{fig:three-levels} and~\ref{fig:genera}.

  \item[ADR-4 --- Pass-through doctrine and enforced full assignment]\hfill\\
  \textbf{Context:} For non-tree-shaped methods (such as the hash cross-check), tree- and
  trie-specific axes like path compression or node type cannot meaningfully be assigned.
  \textbf{Decision:} Such axes are not dropped but assigned as pass-through variants
  (\texttt{None}) --- a pass-through value is an algorithm with an empty path, not an omission;
  implemented is the pass-through building block for path compression
  (\texttt{Path\allowbreak{}Compression\allowbreak{}None}), the node type possesses no None variant
  and is regularly assigned one of the four ART node formats, which a non-tree-shaped building block
  does not use (design; without a code building block); the mandatory concept
  \texttt{IsComposition} of the genus \texttt{SearchAlgorithm} enforces at compile time the
  assignment of all eighteen organs (eighteen type requirements and an organ counter of eighteen
  with a check for ABI break), so that no composition of this genus with an open axis compiles.
  \textbf{Consequence:} Non-tree-shaped methods require no family of their own but a restricted
  axis subset within the same genus.
  \textbf{Evidence:} Sections~\ref{sec:sota-instances},~\ref{sec:anatomy-levels}
  and~\ref{sec:axes-impl}.

  \item[ADR-5 --- Layer model E4--E0 with frozen contracts and contract tests]\hfill\\
  \textbf{Context:} With a product space that rules out any complete materialization, changes
  across several levels at once are practically not acceptable; therefore each level is completed
  individually.
  \textbf{Decision:} The experiment pipeline is organized into the five layers E4--E0 (E4
  definition and evaluation, E3 permutation tree, E2 tier binaries, E1 runtime configuration, E0
  infrastructure) and the cross-cutting concern~M\@; the E identifiers are a conceptual order with
  one code location each. Each layer is completed individually, with an explicitly frozen interface
  contract towards the neighbouring level (the freezing of the contract of layer E1 in the ABI
  adapter: signatures, data layout and capabilities frozen, semantic changes only additive), and
  contract-tested against a compile-time substitute (fake) of the neighbouring level.
  \textbf{Consequence:} Every level can be accepted individually before the next one begins; the
  boundary E3$\to$E2 is contract-tested without a module build against a fake of layer E2, checked
  is solely the static tree and view structure including its mixed-radix bijection.
  \textbf{Evidence:} Sections~\ref{sec:anatomy-layers} and~\ref{sec:impl-contracts};
  Figures~\ref{fig:e-layers} and~\ref{fig:impl-layers}.

  \item[ADR-6 --- Explicitly coded non-observation instead of a phantom value]\hfill\\
  \textbf{Context:} Values that cannot truly be collected (counters without hardware access, axes
  without a real runtime parameter) invite plausible phantom values that would systematically
  poison the heuristic derivation.
  \textbf{Decision:} The system reports explicitly invalid or empty instead of a phantom value and
  distinguishes three states per measured value: valid measured value, explicitly coded
  non-observation (not applicable or source not available) and measurement error\footnote{The code
  vocabulary \texttt{honest-0}/\texttt{honest-100}/\texttt{honest-empty} is the implementation name
  of this principle (Appendix~\ref{app:glossary}); methodically, non-observation is a declared class
  of missing values in the sense of the missing-data literature~\cite{graham2009missing}, never an
  imputed value.}. A validity flag travels with every event, so that a measured zero remains a
  measurement result; the wall-clock source offers only mean and throughput and refuses to label
  its mean as a percentile, and a negative duration is a measurement error, not a zero value; a PMC
  catalogue without a model entry yields \enquote{not applicable}, never a guessed encoding; the
  generators produce no output for a measurement series without a valid row, rather than an empty
  one.
  \textbf{Consequence:} Gaps remain visible and can be closed deliberately; the principle is
  programmed into the building blocks rather than merely documented.
  \textbf{Evidence:} Sections~\ref{sec:measurement-system},~\ref{sec:impl-system-axes}
  and~\ref{sec:axis-triad}.

  \item[ADR-7 --- Four-repository architecture with orthogonal carrier stages]\hfill\\
  \textbf{Context:} Measurement apparatus, probe, experiment definition including evaluation, and
  the manuscript have separate responsibilities and life cycles; the measured values are built
  directly into the manuscript under version control, which is why the manuscript is a repository
  of its own with its own pipeline.
  \textbf{Decision:} Four repositories --- \texttt{comdare-cache-engine} (\emph{how} measuring is
  done; building-block library and builder in the Pitchfork source-tree layout,
  Appendix~\ref{app:glossary}), \texttt{comdare-prt-art} (the probe), the project repository
  \texttt{probst-Diplomarbeit-cache-engine} (\emph{what} is tested by the user: experiment
  definition in XML and XSD, measurement driver, evaluation toolchain) and the manuscript
  repository with its own CI, into which the forward channel deposits the generator output as table
  and diagram files per commit ---; orthogonal to this division lie the four carrier stages planner
  $\to$ CacheEngineBuilder $\to$ tier binary $\to$ hybrid tier binary, where stage $N+1$ hangs only
  on stage $N$; planner and CEB are two binaries --- the planner
  \texttt{comdare-experiment-planner} and the measurement driver \texttt{comdare-messung-driver}
  ---, and the two roles appear as planner line and CEB contract line in the self-stamp of the
  measurement driver (\texttt{--version}) in one driver binary; the measurement aggregator is run
  host-side by the \texttt{ExperimentDriver}.
  \textbf{Consequence:} Builder and CacheEngine share the engine repository --- the builder as an
  executable above its builder library, the CacheEngine as a library ---; the roles remain cleanly
  separated, and the probe is not a stage but the transversal supplier of organ building blocks.
  \textbf{Evidence:} Section~\ref{sec:repos} and Appendix~\ref{app:code};
  Figure~\ref{fig:traeger-stufen}.

  \item[ADR-8 --- Family-wise adapter fallbacks for self-contained runnable builds]\hfill\\
  \textbf{Context:} Third-party code (allocator vendors, SOTA building blocks) is not available
  or activated on every platform and in every build.
  \textbf{Decision:} Compile-time flag families switch adapters on or off --- per vendor one
  family switch, whose use follows from approval and availability, behind 23 vendor shim headers
  ---, each with a safe, family-wise chosen fallback: the allocator adapters replace the missing
  vendor in \texttt{if constexpr} branches with the portable aligned allocation
  (\texttt{portable\_aligned\_alloc}), the vendor shim supplies only forward stubs for the
  never-executed branch, and the axis value \texttt{StdMalloc} is a building block of its own, not
  a stopgap; data-structure adapters of third-party search methods switch via a family switch and
  possess a self-contained re-implementation as their body, while the binding to the original code
  runs via the library build with the original compiler.
  \textbf{Consequence:} The build remains self-containedly runnable without any third-party
  activation.
  \textbf{Evidence:} Section~\ref{sec:axes-impl}.

  \item[ADR-9 --- Measurement gating in three layers and four work modes]\hfill\\
  \textbf{Context:} Measurement and observer couplings cost time and must not burden the
  production path.
  \textbf{Decision:} The measurement couplings of the runtime adapter are switched via the master
  option \texttt{COMDARE\_MEASUREMENT\_MODE}, which defines \texttt{COMDARE\_MEASUREMENT\_ON} and
  which the release mode forcibly switches off; the switchable area falls into three separable
  layers --- base time, observer counters and fine-grained segment timers per axis ---, the
  switch-off stages of the measurement. The run methodology knows four work modes of the same
  permutation, as a registry and not as a label: \texttt{build} (build without measuring),
  \texttt{measure} (deterministic single-thread measurement run, the only mode with measurement),
  \texttt{compare} (reads the measurement store, the stage before \texttt{release}) and
  \texttt{release} (full optimization without measurement instrumentation); the work mode is a
  sub-axis of the measurement realm and does not enter the binary identity, a debug state does not
  exist (\texttt{--debug} is a flag of the planner shell), and the measurement aggregator is run
  host-side by the \texttt{ExperimentDriver}.
  \textbf{Consequence:} The production path stays measurement-free and causes no measurement
  overhead; four modes are selectable, in build semantics only \texttt{measure} differs --- all
  four registry rows carry the release build type.
  \textbf{Implementation status:} The mode-specific replay run of \texttt{compare} and the optimal
  release binary generated from measured values are planned but not implemented;
  \texttt{compare} is a selectable, validatable label with release-equivalent build semantics.
  \textbf{Evidence:} Sections~\ref{sec:repos},~\ref{sec:impl-pipeline} and~\ref{sec:eval-method};
  Figures~\ref{fig:modi-kette} and~\ref{fig:abschaltungsstufen}.

  \item[ADR-10 --- The \texttt{std::map} contract as common denominator with a conformance gate]\hfill\\
  \textbf{Context:} Structurally different search algorithms can only be compared fairly through
  a common, well-understood outer interface.
  \textbf{Decision:} All methods are mapped onto the C++ contract of the ordered map
  \texttt{std::map<Key,\,Value>}~\cite{iso_cpp}; the container genera are held against one
  standard oracle each --- \texttt{Set} against \texttt{std::set}, \texttt{Sequence} against
  \texttt{std::deque}, \texttt{Adapter} against \texttt{std::queue}, \texttt{std::stack} and
  \texttt{std::priority\_queue}, \texttt{View} against \texttt{std::span} ---, and
  \texttt{std::vector} serves as the inner substrate of the genus \texttt{Adapter}. A conformance
  gate checks this contract before every measurement in the binding order import $\to$ gate $\to$
  measure, via the five core operations \texttt{insert}, \texttt{lookup}, \texttt{erase},
  \texttt{clear} and \texttt{size}, with conformance ratios as the result form.
  \textbf{Consequence:} Only comparables are compared; the axes describe exclusively the inner
  behaviour behind the fixed interface.
  \textbf{Implementation status:} The extension of the gate to the full \texttt{std::map}
  interface and the genus iterator library are not implemented; the requirement to run the
  container genera through a complete \texttt{std::vector} hull is not met --- the implemented
  Sequence oracle is \texttt{std::deque}.
  \textbf{Evidence:} Sections~\ref{sec:measurement-system} and~\ref{sec:impl-contracts}.

  \item[ADR-11 --- Design-pattern register with evidenced roles]\hfill\\
  \emph{Command as bundling and command pattern.}\hfill\\
  \textbf{Context:} The apparatus is deliberately a composition of named design patterns; unclear
  pattern roles would dilute the architecture.
  \textbf{Decision:} Used are, each with a named code location: \emph{Strategy} with
  CRTP~\cite{coplien1995crtp} per organ axis (the stamp contract of the carriers is a CRTP contract
  as well, ADR-13), \emph{Abstract Factory} (hybrid docks, composition factory, probe registry),
  \emph{Director} with \emph{Builder} (plan emission), \emph{Iterator} (the lazily traversed
  experiment tree), \emph{Adapter} (third-party code behind the ABI boundary), \emph{Observer}
  (measurement snapshots per axis), \emph{Visitor} (the measurement seam at the \texttt{dlopen}
  transition, ADR-15), \emph{Chain of Responsibility} (hardware collection chain and selection
  filter chain), \emph{Memento} (copy-on-write in the ABI adapter, state per tier), \emph{Command}
  and \emph{Facade} --- the same roles as the pattern register in
  Section~\ref{sec:impl-hybrid-lager} (Figure~\ref{fig:muster}), named after the catalogue of
  Gamma et al.~\cite{gamma1994patterns} and composed as a tuple of axis
  policies~\cite{alexandrescu2001modern}. The patterns \emph{Composite} and \emph{Decorator} are not
  among them: the measurement hulls of the axes hold no component instance and no full interface
  and are expressly not a Decorator, and a Composite in the sense of the catalogue has no
  counterpart in the code.
  \emph{Not} used in the hot path and in the organ axes are
  \emph{Singleton}, \emph{Flyweight} and \emph{Bridge}; host-side registries use Meyers
  singletons\footnote{Meyers singleton: the one instance as a function-local static object that is
  initialized on the first call~\cite{meyers2005effective}.}, the catalogue generation uses
  flyweight tables per axis, and \enquote{Bridge} is in the code only a colloquial designation for
  connecting pieces (ABI adapter, CI bridge jobs, Sandy Bridge), not a pattern. The \emph{Command}
  pattern carries the compile-time bundling of axis recombinations (the system complex axis as
  command wrapper, the hub of the external tools as invoker, the command base of the axes in the
  cross-cutting concern~M) and the typed control command channel; the command channel of the hybrid
  mode is one of these roles.
  \textbf{Consequence:} Every pattern carries named roles with code location in the wiring of the
  apparatus and is executed at compile time rather than runtime-polymorphically.
  \textbf{Evidence:} Sections~\ref{sec:anatomy-patterns},~\ref{sec:heuristics-modes}
  and~\ref{sec:impl-hybrid-lager}.

  \item[ADR-12 --- Dual measurement regime: immutable Talos and root Linux]\hfill\\
  \textbf{Context:} Production fidelity and measurement depth are in conflict: a
  production-faithful system denies precisely the hardware-counter access that deep measurements
  require.
  \textbf{Decision:} On the production target machines, every measurement is collected under two
  operating-system regimes: an immutable operating system (Talos\footnote{Talos Linux: the
  immutable, API-driven operating system for Kubernetes by Sidero
  Labs~\cite{siderolabs_talos}.}) as the production mirror and a root Linux with full
  hardware-counter access (\texttt{perf}~\cite{demelo2010perf,linux_perf_event_open} and the
  model-specific registers, MSR).
  \textbf{Consequence:} The time- and observer-based categories run identically under both
  regimes and secure comparability; the PMC categories are collected only under the privileged
  regime.
  \textbf{Implementation status:} The decision stands; implemented is the root-Linux half: the
  reported measurements all originate from the root-Linux regime of the two production machines
  --- prod1 with the AMD microarchitecture Zen~5 and prod2 with the Intel Core i9-12900K (Alder
  Lake) ---; the Talos regime is planned but not measured. For Intel, the PMC raw-event catalogue
  contains no assignment; for Zen~5 it is cross-checked by measurement.
  \textbf{Evidence:} Sections~\ref{sec:eval-method} and~\ref{sec:contributions}.

  \item[ADR-13 --- Four carrier stages and one stamp contract]\hfill\\
  \emph{The stamp as identity, store key, recognition criterion and classification.}\hfill\\
  \textbf{Context:} The configuration travels through four carrier stages (planner,
  CacheEngineBuilder, tier binary, hybrid tier binary); each stage needs an ABI-stable,
  compile-time identity into which no measured value enters.
  \textbf{Decision:} A single CRTP contract binds all four carriers: seven stamp interfaces
  (version, measurement line, system line, organ line, fingerprint, overall stamp, attached lines)
  times four carriers yield 28 three-valued admission cells (Table~\ref{tab:stempel-matrix}); the
  absence of the organ line at the CEB is the static barrier \enquote{the CEB carries no organ
  axes}, and the check runs as a \texttt{static\_assert} under every heir. The planner carries a
  SHA-256 fingerprint from version, ISA and operating system, the tier binary a SHA-512 fingerprint
  over eleven members (format~6), among them the three realm lines, the toolchain member and the
  hash of the overlay source set (concatenation in fixed order, SHA-512 once); the CEB carries a
  SHA-512 provenance fingerprint, which is not a reuse criterion. The bracket grammar of the stamp
  entries is defined at one place (\texttt{c\{p.e\}} is \texttt{c} with the sub-members \texttt{p}
  and \texttt{e}, never flat), and the bracket depth is the meta-meta level of the entry.
  \textbf{Consequence:} The stamp fulfils five roles --- identity, cache key, store key, recognition
  criterion and classification; the self-stamp of the driver binary (\texttt{--version}) carries
  only the classification and is neither cache key nor build input.
  \textbf{Implementation status:} The registry-wrapper stamp of the SOTA branch and the stamp
  export of the hybrid stage are not implemented; the stamp identity in every diagram and every
  table of the appendix is likewise not implemented.
  \textbf{Evidence:} Sections~\ref{sec:stamp-model},~\ref{sec:anatomy-levels},~\ref{sec:mess-chain},~\ref{sec:impl-contracts}
  and~\ref{sec:eval-identity}; Figure~\ref{fig:traeger-flaechen}.

  \item[ADR-14 --- Identity cut: only organ axes form the \texttt{binary\_id}]\hfill\\
  \emph{System and measurement axes are identity-neutral.}\hfill\\
  \textbf{Context:} Three axis realms (organ, system, measurement) each carry a registry of their
  own with generated XML\@; if measurement or system values entered the identity, duplicate
  binaries of the same permutation would arise, and a store that holds the same binary several
  times.
  \textbf{Decision:} The \texttt{binary\_id} is the path of eighteen \texttt{axis=value} segments
  in canonical slot order (the organ count eighteen is defined at exactly one place); all
  measurement axes and all system axes carry \texttt{binary\_id = never} in their registries;
  sub-axes are runtime permutation and appear in no build-job legend; the work mode is
  identity-neutral. The build world instead travels along as a fingerprint member --- the system
  line and the toolchain member with \texttt{vendoropt} (ADR-17).
  \textbf{Consequence:} Store key and regression detector (reference catalogue $2^{17}$ with
  CRC-64 checksum) are organ-stable; the same organ permutation against two build worlds is a
  different artefact with a different fingerprint.
  \textbf{Evidence:} Sections~\ref{sec:axis-triad},~\ref{sec:stamp-model}
  and~\ref{sec:impl-system-axes}; Table~\ref{tab:axes-overview}.

  \item[ADR-15 --- Measurement-visitor seam (surface 3) instead of a sidecar]\hfill\\
  \emph{Family and genus interfaces remain unchanged.}\hfill\\
  \textbf{Context:} An observer that fetches fields from a tier binary from the outside measures
  the neighbourhood and not the object: what is measured is then decided by the host through the
  choice of fields; measurement fields in the domain interfaces would alter the family and genus
  contracts.
  \textbf{Decision:} Each carrier stage offers three interface surfaces: surface~1 the genus
  interface (runtime, abstract factory), surface~2 the stamp (compile time, ABI-stable), surface~3
  the measurement cut-through; the measurement crosses as a seam of its own
  (\texttt{tier\_measure\_accept} with \texttt{IMessVisitor}), a visitor at the \texttt{dlopen}
  transition with double dispatch~\cite{gamma1994patterns}; depending on whether the measurement
  facility is activated in CEB and tier binary, a measurement visitor is handed over at the
  interface or not.
  \textbf{Consequence:} The domain interfaces contain no measurement fields; the gate combination
  CEB off, tier on is evidenced by a coverage test; the measurement-gate assignment travels along
  as a fingerprint member of its own with positions of its own for the hybrid level.
  \textbf{Evidence:} Sections~\ref{sec:anatomy-levels},~\ref{sec:impl-pipeline}
  and~\ref{sec:measurement-system}; Figure~\ref{fig:traeger-flaechen}.

  \item[ADR-16 --- Store as stock with compile-time keys]\hfill\\
  \emph{Measurement always on, in the store only re-measure what is missing.}\hfill\\
  \textbf{Context:} A reference space of $2^{17}$ compositions per system permutation (according
  to the catalogue header about 224~GB and 24~h of complete build) is manageable only with reuse of
  built binaries and measured values; the measurement is to be always switched on and, in the
  store, to re-measure only what cannot be found there.
  \textbf{Decision:} Stock~1 holds the binary under its SHA-512 fingerprint; stock~2 holds the
  measured values under the pair of fingerprint of the measured binary and hardware identity of the
  measuring machine, the measurement cell stands beside it and is never fused into the digest; a
  binary found is skipped if its passed test attestation is inventoried, otherwise built or
  measured; a version lock over the overlay source set turns unnoticed source changes into a CI
  error.
  \textbf{Consequence:} The identity travels unchanged from the store via the measurement line
  (\texttt{permutation\_id}, \texttt{fingerprint}) into workbook and appendix.
  \textbf{Implementation status:} Stocks~1 and~2 and the key schema are implemented; stocks~3
  and~4 (synthesis batches and solutions), the planner leaf function, the CEB full placement, the
  hybrid storage path, the commit skip, the rebuild flag per stage and the XML section
  \texttt{publish} form the full store build-out and are not implemented.
  \textbf{Evidence:} Sections~\ref{sec:lager},~\ref{sec:eval-method},~\ref{sec:stamp-model}
  and~\ref{sec:heuristics-modes}; Figures~\ref{fig:lager-bestaende}
  and~\ref{fig:eval-lager-kreislauf}.

  \item[ADR-17 --- Optimization level as the user's choice]\hfill\\
  \emph{House-wide O2, O3 on XML request; the build world is identity-effective.}\hfill\\
  \textbf{Context:} Measurement series over differently optimized build worlds are not
  comparable; the maximum optimization must remain selectable for the user.
  \textbf{Decision:} The release flags apply house-wide with O2 --- including vendor build,
  FetchContent and embedding --- via a forced cache value; O3 is built alongside only on XML input
  request; the build-world level enters the toolchain member of the fingerprint as the field
  \texttt{vendoropt} (compile error instead of a silent default); a gate rejects self-supplied
  release optimization flags that bypass the axis; the sub-axis \texttt{opt\_level} offers O0 to
  Ofast with the determinism flag according to IEEE~754~\cite{ieee754_2019}, which only Ofast
  negates.
  \textbf{Consequence:} Two builds of the same permutation against an O2 and an O3 world are
  different artefacts with different fingerprints; a silent skip is ruled out.
  \textbf{Implementation status:} O0 and O1 as the user's choice in the XML and the pass-through of
  further compiler flags from the XML are planned but not implemented (default O2, O3 selectable
  as override).
  \textbf{Evidence:} Sections~\ref{sec:eval-build-scope},~\ref{sec:impl-system-axes}
  and~\ref{sec:complex-axis}.

  \item[ADR-18 --- One official build route as a GNU front end over CMake]\hfill\\
  \emph{\texttt{configure.sh}, \texttt{make}, \texttt{make check} and \texttt{make install}.}\hfill\\
  \textbf{Context:} The only procedure official under Linux is \texttt{configure.sh},
  \texttt{make} and \texttt{make install} or \texttt{make check}, respectively, and these three
  must lie in the root folder; a configuration may have only one source.
  \textbf{Decision:} \texttt{configure.sh} is a POSIX-sh front end without Python, building
  continues to be done with CMake; the \texttt{Makefile} is checked in statically and only
  forwards to CMake, with the GNU standard targets~\cite{gnu2026standards} \texttt{all},
  \texttt{install}, \texttt{check}, \texttt{installcheck}, \texttt{uninstall}, \texttt{clean},
  \texttt{distclean}, \texttt{mostlyclean} and \texttt{maintainer-clean} (the test target is called
  \texttt{check}); the configuration is stored exclusively in \texttt{config.status} and in the
  CMake cache, and the Makefile deliberately contains no default values --- a generated Makefile
  would be a second configuration source without benefit.
  \textbf{Consequence:} There is no second build route beside CMake and no second configuration
  source; the front end is checkable with \texttt{sh -n}.
  \textbf{Implementation status:} \texttt{configure.sh} and \texttt{Makefile} lie in the
  repository; the framework export for consumers (\texttt{install(EXPORT)},
  \texttt{find\_package}) and the follow-up checks are not implemented.
  \textbf{Evidence:} Appendix~\ref{app:code}.

  \item[ADR-19 --- Four CI cells without skip]\hfill\\
  \emph{gcc and clang each as release and debug; a job fails or passes.}\hfill\\
  \textbf{Context:} Two compilers find different errors; a silent job skip or a tolerated error
  feigns a passed run.
  \textbf{Decision:} The test suite runs in four cells --- gcc and clang each as release and debug
  build --- without change skip; the cell is the warning, the job fails; \texttt{allow\_failure}
  carries only two declared exceptions --- the arm64 smoke build without an available runner and
  the manual relock maintenance job ---, no build or test cell; per compiler and build type there
  is one object-cache key with exactly one writer; the minimum compiler level is g++~15.3; the
  golden reference catalogue $2^{17}$ with CRC-64 checksum is the regression detector with frozen
  cardinality.
  \textbf{Consequence:} Accepted is only a run without failed tests in all four cells without an
  exception list --- a gate in the sense of continuous delivery~\cite{humble2010cd} ---, and a
  test count is named only together with the run that produced it.
  \textbf{Implementation status:} The mini-pipeline per carrier stage (four stages times gcc and
  clang times release and debug) is planned but not implemented.
  \textbf{Evidence:} Section~\ref{sec:impl-contracts}.

  \item[ADR-20 --- XML as the only control]\hfill\\
  \emph{The XML carries the what, where and when, never the how.}\hfill\\
  \textbf{Context:} The apparatus is operated from a declarative experiment profile; terms are not
  translated but must be conformant --- were a translation necessary, that is a regression and a
  compile-time error.
  \textbf{Decision:} Two schemas describe profile and measurement series
  (\texttt{experiment\_schema.xsd}, \texttt{messreihe\_v32\_schema.xsd}); the parse seam aborts
  with an error message, the validation is a building block of its own, and the CI checks
  well-formedness; the XML describes the what, where and when, never the how; the axis registries
  are generated as XML and are thus machine-readable images; the planner is a subcommand tool
  (\texttt{validate}, \texttt{plan}, \texttt{cache-key}, \texttt{fingerprint}, \texttt{status},
  \texttt{check-size}, \texttt{simulate}, \texttt{version}, \texttt{help}) and the shell control
  path; as the second control path, an API command line with a headless service is planned, which
  serves the planner as a cell on Kubernetes or on bare metal, as well as a command-line-guided
  creation mode for the user XML in which unavailable options are not selectable.
  \textbf{Consequence:} There is no side route beside the XML\@; the SOTA profiles are executable
  as XML\@.
  \textbf{Implementation status:} The headless control path and the creation mode are not
  implemented; the hybrid section is parsed but not declared in the schema, the organ meta-meta
  section declared but not read, and the \texttt{publish} section is missing.
  \textbf{Evidence:} Sections~\ref{sec:mess-chain} and~\ref{sec:impl-pipeline};
  Figure~\ref{fig:pipeline}.

\end{description}
